\begin{afterword}
\summaryhead

The post names a fallacy: raising one hypothesis for consideration when no evidence singles it out. Its example is a detective who, with no clues to a murder in a town of a million, proposes that one particular resident did it. The author argues that locating a hypothesis in a large space takes most of the evidence needed to believe it. Narrowing a billion possibilities to ten takes 27 bits; 4 more give better than even odds. A hypothesis raised without that evidence has skipped most of the work, and because people treat whatever they think about as ``in the running'', raising it does much of the work of persuasion.

The post applies the rule to theists who take gaps in science as evidence for Jehovah, to creationists who move doubt within evolutionary theory onto intelligent design, and to quantum mechanics. There, the author argues, no evidence favors collapse, and doubt about the Born probabilities should stay among theories with the usual properties of physical law, all of which single-world theories break.

\responsehead

The general rule is sound, and the post states it with care. It asks for some evidence when a hypothesis is first raised, not for proof, and it allows evidence a court would not admit. Its arithmetic is correct. Its claim that thinking about a hypothesis does much of the work of persuasion is offered without evidence, but research supports it: Koehler's 1991 review found that people who explain or imagine a possibility become more confident that it is true. The theist and creationist cases are the familiar ``God of the gaps'' criticism. The post's addition is the demand for evidence that points to one god, or to design, in particular.

The quantum example needs more than the rule. Nothing singles out the detective's suspect, Mortimer, from a million people. In quantum mechanics, the evidence that picked out the theory picked out the collapse and no-collapse readings together, since they predict the same results in every experiment done so far. The charge of privilege therefore depends on a further premise: that a collapse postulate is so complex that its rivals are ``no more complicated or unlikely''. That premise is the simplicity argument for many-worlds, and it is the point in dispute; the book argues for it in ``Many Worlds, One Best Guess''. In the published transcript of the conversation the post reports, Aaronson disputed exactly this premise, saying that one world adds less complexity than the alternatives the author proposed. The post reports Aaronson's concession about evidence and leaves out the reason.

The supporting claims are weak. ``There must be a trillion better ways'' to answer the Born question comes with no example; the one candidate the author had discussed in 2008, Hanson's mangled worlds, was put ``under fifty percent''. The closing sentence says single-world theories violate ``all'' the listed properties of physical law. Bell's theorem does make every single-world account nonlocal. But Bohmian mechanics, a single-world theory, keeps the usual linear evolution of the wave function and moves particles deterministically and continuously.

In short: a sound rule about evidence and attention, with correct arithmetic, applied to quantum mechanics through a premise about simplicity that is the point in dispute, and closed with a claim about all single-world theories that Bohmian mechanics does not fit.
\end{afterword}
